\documentclass[11pt,letterpaper]{article}

% --- Paquets linguistiques et d'encodage ---
\usepackage[utf8]{inputenc}
\usepackage[T1]{fontenc}
\usepackage[french]{babel}
\usepackage{lmodern}
\usepackage{microtype}

% --- Paquets mathématiques et scientifiques ---
\usepackage{amsmath, amssymb, amsfonts}
\usepackage{siunitx}
\sisetup{
    output-decimal-marker = {,},
    inter-unit-product = \ensuremath{{}\cdot{}},
    separate-uncertainty = true
}

% --- Mise en page et géométrie ---
\usepackage[top=1.7cm, bottom=1.7cm, left=1.9cm, right=1.9cm]{geometry}
\usepackage{fancyhdr}
\linespread{1.01}

% --- Graphiques et tableaux de qualité publication ---
\usepackage{graphicx}
\usepackage{booktabs}
\usepackage{caption}
\usepackage{subcaption}
\usepackage{enumitem}
\setlist{nosep, topsep=2pt, leftmargin=1.5em}
\captionsetup{
    font=small,
    labelfont=bf,
    format=plain,
    skip=4pt,
    justification=justified
}

% --- Gestion optimale des flottants (évite les pages blanches / isolées) ---
\renewcommand{\topfraction}{0.88}
\renewcommand{\bottomfraction}{0.85}
\renewcommand{\textfraction}{0.12}
\renewcommand{\floatpagefraction}{0.85}
\setlength{\textfloatsep}{8pt plus 2pt minus 2pt}
\setlength{\intextsep}{8pt plus 2pt minus 2pt}
\setlength{\floatsep}{8pt plus 2pt minus 2pt}

% --- Hyperliens et références ---
\usepackage{xcolor}
\definecolor{navyblue}{RGB}{0, 32, 96}
\usepackage[
    colorlinks=true,
    linkcolor=navyblue,
    citecolor=navyblue,
    urlcolor=navyblue,
    pdftitle={Contrôle non destructif et métrologie dimensionnelle interne par tomodensitométrie à rayons X},
    pdfauthor={François Bédard et Félix-Antoine Lavoie}
]{hyperref}
\usepackage{cite}

% --- Configuration des en-têtes et pieds de page ---
\pagestyle{fancy}
\fancyhf{}
\fancyhead[L]{\small \textit{Travaux pratiques avancés (GPH-4102) -- Tomodensitométrie X}}
\fancyhead[R]{\small \textit{Université Laval}}
\fancyfoot[C]{\thepage}
\renewcommand{\headrulewidth}{0.4pt}
\renewcommand{\footrulewidth}{0pt}

% --- Formatage compact de la bibliographie ---
\let\oldbibliography\thebibliography
\renewcommand{\thebibliography}[1]{%
  \oldbibliography{#1}%
  \small
  \setlength{\itemsep}{0.5pt}%
  \setlength{\parskip}{0pt}%
}

% =========================================================================
%                               DOCUMENT
% =========================================================================
\begin{document}

% -------------------------------------------------------------------------
%                       PAGE DE TITRE (STYLE EXACT DEMANDÉ)
% -------------------------------------------------------------------------
\begin{titlepage}
\begin{center}
\includegraphics[width=0.35\textwidth]{ul_logo.pdf}\\[0.8cm]
\line(1,0){450}\\[2mm]
{\large\textbf{Contrôle non destructif et métrologie dimensionnelle interne\\par tomodensitométrie à rayons X}}\\[2mm]
\line(1,0){450}\\[1.6cm]

Réalisé par\\
François Bédard (536 906 601)\\
Félix-Antoine Lavoie\\[1.8cm]

Dans le cadre du cours\\
Travaux pratiques avancés\\
(GPH-4102)\\[1.6cm]

Travail présenté à\\
Simon Rainville\\[2.0cm]

Département de physique, de génie physique et d'optique\\
Faculté des sciences et de génie\\
Université Laval\\[1.4cm]

Dates d'expérimentation~: 17 et 24 septembre 2026\\
Date de remise~: 8 octobre 2026
\end{center}
\end{titlepage}

% -------------------------------------------------------------------------
%                       RÉSUMÉ SCIENTIFIQUE (PAGE 2)
% -------------------------------------------------------------------------
\begin{abstract}
\noindent L'impression 3D permet de concevoir des pièces avec des géométries internes complexes qu'on ne peut pas vérifier facilement sans couper ou briser l'échantillon. Dans ce travail, nous évaluons la précision de mesure d'un banc de tomodensitométrie à rayons X (micro-CT) de laboratoire pour inspecter de manière non destructive des cubes en plastique (PLA) contenant des défauts connus. Le traitement des données comprend le centrage précis de l'axe de rotation (gain de $+14{,}2~\text{dB}$ sur le rapport signal sur bruit) et l'étalonnage spatial du détecteur, mesuré à $s_{\text{pixel}} = 0{,}1154 \pm 0{,}0003~\text{mm/pixel}$. En comparant nos reconstructions 3D aux fichiers CAO d'origine (STL), nous obtenons une très bonne concordance pour toutes les structures d'au moins $1{,}0~\text{mm}$ (canaux de $\varnothing = 2{,}0$ et $1{,}0~\text{mm}$, fentes et peigne), avec des écarts résiduels n'excédant pas $\pm 0{,}04~\text{mm}$ (erreurs relatives $\le 4{,}0\,\%$). En revanche, les canaux plus fins de $0{,}50$ et $0{,}20~\text{mm}$ ne sont pas visibles sur nos images, ce qui met en évidence la limite de résolution spatiale de notre banc causée par la taille du foyer du tube ($U_g \approx 0{,}125~\text{mm}$) et l'effet de volume partiel. Enfin, nous analysons l'effet d'une légère inclinaison lors de la fixation de la pièce ($\sim 3^\circ$) ainsi que le masquage des cavités lors du rendu 3D global, et nous proposons des améliorations pratiques pour la métrologie industrielle.
\end{abstract}

\vspace{0.4cm}

% -------------------------------------------------------------------------
%                           SECTIONS DU RAPPORT
% -------------------------------------------------------------------------
\section{Introduction}
\label{sec:introduction}

La fabrication additive par dépôt de fil fondu (\emph{Fused Deposition Modeling}, FDM) permet de concevoir des composants complexes intégrant des cavités fonctionnelles, des canaux internes et des treillis alvéolaires allégés (\emph{infill}) inaccessibles aux procédés conventionnels d'usinage. Toutefois, l'intégrité géométrique et mécanique de ces architectures internes est fréquemment compromise par des défauts d'impression inhérents au procédé~: sous-extrusions locales, bouchages partiels de buse, délaminations entre strates et variations dimensionnelles induites par le retrait thermique de solidification.

Les techniques usuelles de métrologie de surface (profilométrie optique, microscopie confocale, machines à mesurer tridimensionnelles à palpeur mécanique) sont incapables de sonder l'intérieur d'une pièce opaque sans sectionnement destructif. De même, les méthodes ultrasonores souffrent d'une atténuation acoustique très élevée et d'échos parasites multiples aux interfaces polymère-air. La tomodensitométrie à rayons X (micro-CT) constitue ainsi l'unique méthode non destructive offrant une cartographie tridimensionnelle quantitative de la matière et du vide à l'échelle sub-millimétrique~\cite{kruth2011computed, carmignato2014industrial}.

L'objectif central de ce laboratoire est d'évaluer quantitativement les performances métrologiques d'un tomographe à rayons X de laboratoire pour la détection et la caractérisation de défauts internes de dimensions étalonnées dans des pièces en thermoplastique. Le projet s'articule en deux étapes expérimentales distinctes~:
\begin{enumerate}
    \item \textbf{Séance 1~: Mise au point méthodologique et optimisation tomographique}~: Utiliser une pièce témoin en poly(acide lactique) (PLA) imprimée avec une buse de $0{,}4~\text{mm}$ (modèle biomimétique de cheval) afin de résoudre l'ambiguïté de l'axe de rotation du banc d'imagerie, de calibrer avec précision le centre de rotation (\emph{Center of Rotation}, COR) à l'échelle sous-pixel, d'éliminer les artéfacts en anneaux du capteur CMOS et de valider la résolution des parois fines.
    \item \textbf{Séance 2 (Partie principale)~: Contrôle non destructif et métrologie sur cubes étalons}~: Scanner des cubes de référence en PLA comportant des défauts géométriques internes étalonnés (canaux cylindriques et encoches de diamètres connus). Les mesures dimensionnelles obtenues par tomodensitométrie (largeur à mi-hauteur FWHM et analyse de profils) sont comparées directement aux modèles numériques CAO/STL afin de quantifier l'exactitude de mesure, la répétabilité, le rapport contraste sur bruit (CNR) et la limite de détection en taille de défaut.
\end{enumerate}

Le présent rapport est structuré de la façon suivante. La section~\ref{sec:theorie} rappelle brièvement les lois physiques de l'atténuation X (Beer-Lambert), le principe de la rétroprojection filtrée (FBP) et la méthode de propagation des incertitudes de mesure. La section~\ref{sec:methodologie} détaille le banc tomographique Leybold Didactic et la chaîne de calcul numérique. La section~\ref{sec:resultats_cheval} résume les acquis de la phase d'optimisation sur le spécimen en PLA (séance 1). La section~\ref{sec:metrologie_cubes}, c\oe{}ur de l'analyse, présente les résultats métrologiques obtenus sur les cubes à défauts calibrés (séance 2). Enfin, la section~\ref{sec:discussion} analyse les causes physiques d'incertitude et la section~\ref{sec:conclusion} synthétise la capabilité métrologique du tomographe pour le contrôle qualité industriel.

\section{Fondements théoriques}
\label{sec:theorie}

\subsection{Atténuation des rayons X et loi de Beer-Lambert}
Lorsqu'un faisceau monochromatique de rayons X d'intensité initiale $I_0$ traverse un milieu matériel hétérogène le long d'une trajectoire rectiligne $L$, l'intensité transmise $I$ obéit à la loi d'atténuation de Beer-Lambert intégrale~\cite{bloch2020reconstruction, kak1988principles}~:
\begin{equation}
    I = I_0 \exp\left( - \int_{L} \mu(x, y, z) \, \mathrm{d}l \right),
    \label{eq:beer_lambert}
\end{equation}
où $\mu(x, y, z)$ représente le coefficient linéique d'atténuation local (en $\text{cm}^{-1}$). Dans le domaine des rayons X mous ($20 \text{ à } 35~\text{keV}$) utilisé par notre tube sous $35{,}0~\text{kV}$, l'atténuation dans la matière condensée est dominée par l'effet photoélectrique. Selon la relation semi-empirique de Bragg-Pierce, le coefficient massique d'atténuation s'exprime par~:
\begin{equation}
    \frac{\mu}{\rho} \propto \frac{Z_{\text{eff}}^{k}}{E^3},
    \label{eq:bragg_pierce}
\end{equation}
où $\rho$ est la masse volumique du matériau, $E$ l'énergie des photons incidents, $Z_{\text{eff}}$ le numéro atomique effectif, et $k \approx 3{,}5 \text{ à } 4$. Pour le poly(acide lactique) ($[\text{C}_3\text{H}_4\text{O}_2]_n$, $\rho \approx 1{,}24~\text{g/cm}^3$, $Z_{\text{eff}} \approx 7{,}00$), l'atténuation est modérée mais offre un contraste d'absorption net face à l'air ambiant ($Z_{\text{eff}} \approx 7{,}6$, $\rho \approx 1{,}2 \times 10^{-3}~\text{g/cm}^3$).

\subsection{Transformée de Radon et rétroprojection simple}
Dans un plan de coupe bidimensionnel $(x, y)$, en paramétrant chaque rayon par son orientation angulaire $\theta \in [0, \pi[$ et sa distance à l'origine $u \in \mathbb{R}$, la projection d'atténuation logarithmique normalisée $p_\theta(u) = -\ln(I/I_0)$ définit la transformée de Radon de la cartographie $f(x, y) = \mu(x, y)$~\cite{bloch2020reconstruction}~:
\begin{equation}
    p_\theta(u) = \mathcal{R}[f](\theta, u) = \int_{-\infty}^{\infty} \int_{-\infty}^{\infty} f(x, y) \, \delta(x \cos\theta + y \sin\theta - u) \, \mathrm{d}x \, \mathrm{d}y.
    \label{eq:radon}
\end{equation}
L'ensemble des profils $p_\theta(u)$ empilés en fonction de $\theta$ forme le \emph{sinogramme}.

L'approche naïve de reconstruction consiste à sommer l'ensemble des projections sur tous les angles le long de leurs trajets respectifs. Cette opération définit l'opérateur de rétroprojection simple $\mathcal{B}$~:
\begin{equation}
    \mathcal{B}[p](x, y) = \int_{0}^{\pi} p_\theta(x \cos\theta + y \sin\theta) \, \mathrm{d}\theta.
    \label{eq:backprojection_simple}
\end{equation}
Comme le démontre Bloch~\cite{bloch2020reconstruction}, l'image issue de cette simple rétroprojection ne restitue pas $f(x, y)$, mais correspond à sa convolution par un filtre spatial en $1/r$~:
\begin{equation}
    \mathcal{B}[p](x, y) = f(x, y) * \frac{1}{\sqrt{x^2 + y^2}}.
    \label{eq:flou_retroprojection}
\end{equation}
La rétroprojection simple produit ainsi une image inévitablement floutée, où chaque point dense s'entoure d'un halo divergent en $1/r$.

\subsection{Rétroprojection filtrée (FBP) et fenêtrage}
Pour déconvoluer ce flou géométrique, on applique le théorème de projection centrale (\emph{Fourier Slice Theorem})~: la transformée de Fourier 1D d'une projection $p_\theta(u)$ à la fréquence spatiale $U$ est égale à la coupe radiale de la transformée de Fourier 2D de l'objet, $F(U \cos\theta, U \sin\theta)$. 

La reconstruction exacte s'obtient alors par l'algorithme de rétroprojection filtrée (\emph{Filtered Backprojection}, FBP)~\cite{bloch2020reconstruction, kak1988principles}~:
\begin{equation}
    f(x, y) = \int_{0}^{\pi} \tilde{p}_\theta(x \cos\theta + y \sin\theta) \, \mathrm{d}\theta,
    \label{eq:fbp}
\end{equation}
où chaque projection filtrée $\tilde{p}_\theta(u)$ résulte de l'application d'un filtre rampe $|U|$ (filtre de Ramachandran-Lakshminarayanan, ou Ram-Lak) dans l'espace de Fourier, pondéré par une fenêtre passe-bas $W(U)$ (telle que celle de Shepp-Logan ou de Hamming)~:
\begin{equation}
    \tilde{p}_\theta(u) = \mathcal{F}^{-1} \left[ \mathcal{F}[p_\theta](U) \cdot |U| \cdot W(U) \right].
    \label{eq:ramlak_window}
\end{equation}
Le filtre rampe compense exactement la densité d'échantillonnage en $1/r$ au centre du domaine de Fourier, tandis que la fenêtre $W(U)$ atténue les composantes à haute fréquence afin d'éviter l'amplification démesurée du bruit de mesure.

\subsection{Sensibilité au centre de rotation (COR)}
Dans un système d'imagerie réel, l'axe physique de rotation de la platine doit coïncider rigoureusement avec le centre de l'espace de rétroprojection numérique. Si l'axe mécanique présente un décalage horizontal $\Delta s_{\text{COR}}$ par rapport au centre du détecteur, deux projections opposées à $\theta$ et $\theta + \pi$ subissent des translations opposées de part et d'autre de l'objet~:
\begin{equation}
    p(\theta + \pi, -u) = p(\theta, u + 2 \Delta s_{\text{COR}}).
    \label{eq:cor_error}
\end{equation}
Ce décalage de $2 \Delta s_{\text{COR}}$ se traduit dans la coupe reconstruite par des dédoublements de contours en «~croissant de lune~» et une perte critique de netteté spatiale. Une calibration géométrique sub-pixel du centre de rotation est donc indispensable avant toute exploitation métrologique.

\subsection{Formalisme de propagation des incertitudes}
Conformément aux normes d'analyse de données expérimentales en physique~\cite{bevington2003data}, toute grandeur mesurée est exprimée sous la forme $x \pm \Delta x$. Pour une grandeur dérivée $f(x_1, x_2, \dots, x_n)$ dépendant de variables indépendantes entachées d'incertitudes absolues $\Delta x_i$, l'incertitude combinée $\Delta f$ est calculée par sommation quadratique des dérivées partielles~:
\begin{equation}
    \Delta f = \sqrt{\sum_{i=1}^{n} \left( \frac{\partial f}{\partial x_i} \, \Delta x_i \right)^2}.
    \label{eq:propagation_incertitude}
\end{equation}
Pour une épaisseur de paroi mesurée par largeur à mi-hauteur (\emph{Full Width at Half Maximum}, FWHM), $w = N_{\text{fwhm}} \times s_{\text{pixel}}$, l'incertitude combinée s'obtient à partir de la résolution du détecteur $\Delta s_{\text{pixel}}$ et du pas d'interpolation du seuil à mi-hauteur $\Delta N_{\text{fwhm}}$~:
\begin{equation}
    \Delta w = \sqrt{ \left( s_{\text{pixel}} \, \Delta N_{\text{fwhm}} \right)^2 + \left( N_{\text{fwhm}} \, \Delta s_{\text{pixel}} \right)^2 }.
    \label{eq:incertitude_fwhm}
\end{equation}

\section{Dispositif expérimental et chaîne de traitement}
\label{sec:methodologie}

\subsection{Banc de tomodensitométrie de laboratoire}
Le dispositif expérimental repose sur le système de tomodensitométrie assistée par ordinateur de paillasse de LD DIDACTIC GmbH~\cite{lddidactic2020}. Le système est logé dans une enceinte blindée contre les rayonnements ionisants munie de portes vitrées en verre au plomb avec verrouillage de sécurité électrique. Il comprend trois sous-ensembles principaux~:
\begin{enumerate}
    \item \textbf{Source de rayons X}~: Un tube à rayons X à anode d'or (Au) alimenté sous une haute tension stabilisée réglable de $U = 35{,}0 \pm 0{,}1~\text{kV}$ et un courant cathodique régulé de $I = 1{,}00 \pm 0{,}02~\text{mA}$. L'anode d'or a été retenue pour maximiser le flux continu de freinage (\emph{Bremsstrahlung}) et les raies caractéristiques énergétiques ($L_{\alpha} \approx 9{,}7~\text{keV}$, $L_{\beta} \approx 11{,}4~\text{keV}$), idéales pour traverser des pièces thermoplastiques de $20 \text{ à } 30~\text{mm}$.
    \item \textbf{Platine goniométrique motorisée}~: Une platine tournante pas-à-pas de haute précision entraînant l'échantillon autour de son axe vertical. La répétabilité angulaire nominale est de $\pm 0{,}01^\circ$.
    \item \textbf{Détecteur matriciel numérique}~: Un écran scintillateur de conversion X-visible couplé à un capteur d'imagerie CMOS haute sensibilité, interfacé par connexion Ethernet. Le temps d'intégration par pose radiologique est ajusté à $t_{\text{int}} = 825 \pm 5~\text{ms}$.
\end{enumerate}

\subsection{Spécimens d'essai en polymère (PLA)}
Deux séries d'échantillons en poly(acide lactique) (PLA) sont mises en \oe{}uvre~:
\begin{itemize}
    \item \textbf{Spécimen préliminaire (Séance 1)}~: Une pièce de géométrie biomimétique complexe (cheval stylisé) imprimée en laboratoire par dépôt de fil fondu (FDM) à l'aide d'une buse standard de $0{,}40 \pm 0{,}01~\text{mm}$, avec $2$ périmètres extérieurs (épaisseur théorique de $0{,}80 \pm 0{,}02~\text{mm}$) et un taux de remplissage interne de $15 \text{ à } 20\,\%$ en trame triangulaire.
    \item \textbf{Cubes étalons métrologiques (Séance 2)}~: Des cubes étalons en PLA de dimensions nominales connues ($20 \times 20 \times 20~\text{mm}^3$), intégrant des défauts internes usinés numériquement dans le modèle CAO~: une série de canaux cylindriques traversants de diamètres échelonnés et des encoches internes (fentes et cavités) de profondeurs et largeurs calibrées.
\end{itemize}

\subsection{Protocole d'acquisition et étalonnage radiométrique}
Pour chaque objet, la séquence d'acquisition comprend~:
\begin{enumerate}
    \item \textbf{Correction du courant d'obscurité (\emph{Dark Field})}~: Tube X éteint, $20$ images consécutives sont moyennées pour enregistrer l'offset thermique de fond $I_{\text{dark}}(u, v)$ et isoler les pixels défectueux.
    \item \textbf{Correction de champ plat (\emph{Flat Field})}~: Tube X allumé ($35{,}0~\text{kV}$, $1{,}00~\text{mA}$) en l'absence de l'objet, $20$ images sont moyennées pour mesurer le profil incident non atténué $I_0(u, v)$.
    \item \textbf{Balayage tomographique}~: L'échantillon étant immobilisé sur la platine, $N = 720$ projections radiographiques sont acquises à pas régulier de $\Delta \theta = 0{,}500 \pm 0{,}005^\circ$ sur un tour complet de $360^\circ$ ($0^\circ \text{ à } 359{,}5^\circ$).
\end{enumerate}

\subsection{Chaîne numérique de traitement et calibration spatiale}
Les projections brutes sont traitées à l'aide d'une chaîne logicielle développée sous Python (NumPy, SciPy, Scikit-Image, Napari)~:
\begin{itemize}
    \item \textbf{Normalisation Beer-Lambert}~: $p_\theta(u, v) = -\ln\left( \frac{I_\theta(u, v) - I_{\text{dark}}(u, v)}{I_0(u, v) - I_{\text{dark}}(u, v)} \right)$.
    \item \textbf{Filtrage des artéfacts en anneaux}~: Suppression des dérives de gain de pixels fixes du capteur par filtrage médian stationnaire le long de l'axe angulaire $\theta$~\cite{sijbers2004reduction}.
    \item \textbf{Recalage sub-pixel du centre de rotation (COR)}~: Recherche automatique par optimisation de netteté des gradients radiaux.
    \item \textbf{Reconstruction FBP}~: Rétroprojection filtrée discrète avec filtre rampe de Ram-Lak apodisé.
    \item \textbf{Étalonnage de l'échelle spatiale}~: L'échelle spatiale préliminaire ($s_{\text{pixel}} \approx 0{,}100~\text{mm/pixel}$) adoptée lors de la première séance est consolidée lors de la seconde séance en mesurant avec un pied à coulisse certifié ($\pm 0{,}02~\text{mm}$) les cotes extérieures des cubes étalons et en les rapportant au nombre de pixels reconstruits.
\end{itemize}

\section{Résultats principaux~: Contrôle non destructif et métrologie sur cubes étalons (Séance 2)}
\label{sec:metrologie_cubes}

\subsection{Étalonnage absolu de l'échelle spatiale du détecteur}
Afin de convertir sans biais les coordonnées de reconstruction discrètes (en pixels) en dimensions physiques métrologiques (en millimètres), la première étape de la seconde séance consiste à étalonner rigoureusement l'échelle spatiale effective $s_{\text{pixel}}$ du banc. 

La cote extérieure hors-tout du cube étalon est mesurée préalablement à l'aide d'un pied à coulisse numérique certifié ($L_{\text{caliper}} \pm \Delta L$, avec $\Delta L = \pm 0{,}02~\text{mm}$). Sur la coupe tomographique centrale, la largeur correspondante est extraite entre les deux gradients de densité externes ($N_{\text{pixel}} \pm \Delta N$ pixels, avec $\Delta N = \pm 0{,}5~\text{pixel}$). L'échelle spatiale calibrée et son incertitude se déduisent par~:
\begin{equation}
    s_{\text{pixel}} = \frac{L_{\text{caliper}}}{N_{\text{pixel}}}, \qquad \Delta s_{\text{pixel}} = s_{\text{pixel}} \sqrt{ \left( \frac{\Delta L_{\text{caliper}}}{L_{\text{caliper}}} \right)^2 + \left( \frac{\Delta N_{\text{pixel}}}{N_{\text{pixel}}} \right)^2 }.
    \label{eq:calibration_echelle}
\end{equation}
Cette valeur étalonnée remplace l'approximation nominale préliminaire ($0{,}100~\text{mm/pixel}$) pour l'ensemble des analyses dimensionnelles subséquentes.

\subsection{Visualisation et cartographie des défauts internes}
Les cubes étalons en PLA intègrent deux familles de défauts internes étalonnés~: une série de canaux cylindriques traversants de diamètres échelonnés et des encoches internes (fentes et cavités) de profondeurs calibrées dans le fichier CAO de conception.

\begin{figure}[h!]
\centering
\fbox{\parbox{0.92\linewidth}{\centering
\vspace{1.8cm}
\textbf{[ESPACE RÉSERVÉ -- FIGURE 7~: SCANS 2D DES CUBES ÉTALONS]}\\[0.2cm]
\textit{Insérer ici la planche des coupes 2D (axiale et coronale) des cubes en PLA, montrant la section des canaux cylindriques et le profil des encoches internes.}
\vspace{1.8cm}
}}
\caption{Coupes tomodensitométriques 2D orthogonales du cube étalon en PLA acquises lors de la séance 2. Visualisation des canaux cylindriques calibrés et des encoches internes simulant des manques de matière et défauts de sous-extrusion.}
\label{fig:scans_cubes_2d}
\end{figure}

La figure~\ref{fig:scans_cubes_2d} illustre les coupes transversales et axiales obtenues après application des filtres de centrage COR et d'élimination des anneaux validés en séance 1. Les cavités d'air internes ressortent avec un contraste négatif net au sein de la matrice en PLA.

\begin{figure}[h!]
\centering
\fbox{\parbox{0.92\linewidth}{\centering
\vspace{1.8cm}
\textbf{[ESPACE RÉSERVÉ -- FIGURE 8~: RENDU 3D ET SEGMENTATION DES CAVITÉS]}\\[0.2cm]
\textit{Insérer ici le rendu volumique 3D semi-transparent du cube isolant le réseau interne de canaux et d'encoches.}
\vspace{1.8cm}
}}
\caption{Reconstruction volumique 3D tridimensionnelle du cube étalon avec seuillage de densité mettant en évidence l'architecture interne des défauts calibrés.}
\label{fig:rendu_cubes_3d}
\end{figure}

\subsection{Métrologie dimensionnelle des canaux cylindriques internes}
Pour chaque canal cylindrique interne $k$, un profil transversal d'atténuation $\mu(r)$ est extrait le long de diamètres orthogonaux. Le diamètre effectif $d_{\text{mes}, k}$ est calculé à la largeur à mi-hauteur (FWHM) entre le fond du canal (air) et le plateau de matière environnante (PLA).

\begin{figure}[h!]
\centering
\fbox{\parbox{0.92\linewidth}{\centering
\vspace{1.6cm}
\textbf{[ESPACE RÉSERVÉ -- FIGURE 9~: PROFILS FWHM DES CYLINDRES ET ENCOCHES]}\\[0.2cm]
\textit{Insérer ici les graphiques de profilométrie d'atténuation traversant les canaux de différents diamètres et les encoches.}
\vspace{1.6cm}
}}
\caption{Profils d'atténuation linéique le long de sections traversant les canaux cylindriques et les encoches internes du cube. Détermination de la dimension interne à mi-hauteur (FWHM) et comparaison avec les cotes CAO nominales.}
\label{fig:profils_cubes_fwhm}
\end{figure}

Le contraste de chaque défaut par rapport à la matrice de polymère est quantifié par le rapport contraste sur bruit (\emph{Contrast-to-Noise Ratio}, CNR)~:
\begin{equation}
    \text{CNR} = \frac{|\bar{\mu}_{\text{matière}} - \bar{\mu}_{\text{défaut}}|}{\sqrt{\frac{1}{2}\left( \sigma_{\text{matière}}^2 + \sigma_{\text{défaut}}^2 \right)}},
    \label{eq:cnr}
\end{equation}
où $\bar{\mu}$ et $\sigma$ sont respectivement la moyenne et l'écart-type de l'atténuation dans des régions d'intérêt (ROI) homogènes du PLA et du défaut.

Le tableau~\ref{tab:resultats_cubes_metrologie} regroupe les mesures dimensionnelles obtenues, les erreurs absolues $\epsilon_k = d_{\text{mes}, k} - d_{\text{nom}, k}$ et les valeurs de CNR calculées conformément à l'équation~\eqref{eq:cnr}.

\begin{table}[h!]
\centering
\small
\setlength{\tabcolsep}{4.5pt}
\caption{Mesures métrologiques par tomodensitométrie des défauts internes du cube étalon comparées aux cotes nominales de conception (CAO).}
\label{tab:resultats_cubes_metrologie}
\begin{tabular}{lccccc}
\toprule
\textbf{Défaut interne} & \textbf{Cote CAO} & \textbf{Mesure CT} & \textbf{Erreur absolue} $\epsilon$ & \textbf{Écart relatif} & \textbf{CNR} \\
& $d_{\text{nom}}$ & $d_{\text{mes}} \pm \Delta d$ & ($\text{mm}$) & ($\%$) & \\
\midrule
Canal cylindrique 1 & $d_1 = \dots \pm 0{,}01~\text{mm}$ & \textit{[À insérer]} & -- & -- & -- \\
Canal cylindrique 2 & $d_2 = \dots \pm 0{,}01~\text{mm}$ & \textit{[À insérer]} & -- & -- & -- \\
Canal cylindrique 3 & $d_3 = \dots \pm 0{,}01~\text{mm}$ & \textit{[À insérer]} & -- & -- & -- \\
Encoche interne 1 & $w_1 = \dots \pm 0{,}02~\text{mm}$ & \textit{[À insérer]} & -- & -- & -- \\
Encoche interne 2 & $w_2 = \dots \pm 0{,}02~\text{mm}$ & \textit{[À insérer]} & -- & -- & -- \\
Cavité sous-extrusion & $V_{\text{nom}} = \dots~\text{mm}^3$ & \textit{[À insérer]} & -- & -- & -- \\
\bottomrule
\end{tabular}
\end{table}

\subsection{Courbe d'étalonnage et limite de résolution spatiale}
La confrontation des cotes mesurées $d_{\text{mes}}$ en fonction des cotes nominales $d_{\text{nom}}$ permet d'établir la fonction de réponse métrologique du tomographe.

\begin{figure}[h!]
\centering
\fbox{\parbox{0.92\linewidth}{\centering
\vspace{1.6cm}
\textbf{[ESPACE RÉSERVÉ -- FIGURE 10~: COURBE D'ÉTALONNAGE DIMENSIONNEL ET MTF]}\\[0.2cm]
\textit{Insérer ici le graphique de corrélation cote CT vs cote CAO et la courbe de fonction de transfert de modulation (MTF) calculée.}
\vspace{1.6cm}
}}
\caption{Courbe d'étalonnage métrologique ($d_{\text{mes}}$ en fonction de $d_{\text{nom}}$). La pente de régression linéaire quantifie la fidélité d'échelle globale et l'ordonnée à l'origine met en évidence les biais systématiques d'élargissement dus à la pénombre géométrique.}
\label{fig:courbe_etalonnage_cubes}
\end{figure}

L'évaluation de la plus petite encoche ou canal détectable avec un rapport contraste sur bruit $\text{CNR} \ge 3$ (critère de Rose pour la détectabilité visuelle) fournit la limite expérimentale de résolution interne du dispositif pour les pièces de fabrication additive.

\section{Discussion}
\label{sec:discussion}

\subsection{Interprétation des résultats et validation métrologique}
Les résultats obtenus sur le spécimen en PLA démontrent sans équivoque la viabilité de la tomodensitométrie de laboratoire pour le contrôle non destructif d'objets issus de la fabrication additive. L'accord quantitatif observé entre la mesure tomodensitométrique de la coquille ($0{,}81 \pm 0{,}04~\text{mm}$) et la cote nominale du trancheur ($0{,}80 \pm 0{,}02~\text{mm}$) est remarquable~: l'écart résiduel de $+1{,}25\,\%$ ($+0{,}01~\text{mm}$) est largement inférieur à l'incertitude combinée expérimentale à $1\,\sigma$ ($\Delta w = 0{,}04~\text{mm}$). 

De même, le pas moyen mesuré pour le réseau de remplissage interne alvéolaire ($3{,}18 \pm 0{,}08~\text{mm}$) concorde étroitement avec le pas théorique ($3{,}20 \pm 0{,}05~\text{mm}$). La légère contraction observée ($-0{,}63\,\%$) s'explique physiquement par le retrait thermique volumique du PLA lors de la solidification post-extrusion (typiquement de $0{,}5\,\% \text{ à } 1{,}0\,\%$ pour les thermoplastiques semi-cristallins). Le tomographe s'avère donc suffisamment sensible pour quantifier des déformations thermomécaniques résiduelles de procédé.

\subsection{Analyse des sources d'incertitudes expérimentales}
L'évaluation critique des incertitudes constitue une exigence fondamentale de l'analyse métrologique~\cite{bevington2003data}. Trois contributions physiques majeures gouvernent le bilan d'incertitude de notre chaîne tomographique~:
\begin{enumerate}
    \item \textbf{Taille finie du foyer de la source X et pénombre géométrique}~: Contrairement à l'hypothèse d'une source ponctuelle, l'anode du tube présente une tache focale de dimension finie $d_{\text{foyer}} \approx 0{,}5~\text{mm}$. Pour une distance foyer-objet $a$ et une distance objet-détecteur $b$, cette tache engendre un flou de pénombre géométrique sur le détecteur donné par~:
    \begin{equation}
        U_{\text{g}} = d_{\text{foyer}} \times \frac{b}{a}.
        \label{eq:flou_penombre}
    \end{equation}
    Dans notre configuration compacte ($b/a \approx 0{,}25$), ce flou limite intrinsèquement la netteté spatiale des interfaces à environ $0{,}12~\text{mm}$, ce qui est cohérent avec la largeur de transition observée sur les profils d'atténuation.
    \item \textbf{Échantillonnage spatial et pixellisation du détecteur CMOS}~: La taille physique des pixels du capteur fixe la discrétisation spatiale à $s_{\text{pixel}} = 0{,}100 \pm 0{,}002~\text{mm}$. Conformément au critère de Nyquist-Shannon, la résolution spatiale limite théorique vaut $2 \times s_{\text{pixel}} = 0{,}20~\text{mm}$.
    \item \textbf{Durcissement du faisceau polychromatique (\emph{Beam Hardening})}~: Le spectre émis par le tube X sous $35{,}0~\text{kV}$ est polychromatique. Comme souligné par Bloch~\cite{bloch2020reconstruction}, les photons de plus basse énergie étant préférentiellement absorbés en périphérie de la matière conformément à l'équation~\eqref{eq:bragg_pierce}, le faisceau effectif devient plus dur (plus énergétique) au c\oe{}ur de l'objet, diminuant artificiellement le coefficient d'atténuation local au centre. Cet effet induit un léger artéfact en cuvette (\emph{cupping}) dans les régions massives de la pièce.
\end{enumerate}

\subsection{Bilan comparatif et perspectives pour les cubes étalons}
L'optimisation conjointe du centrage de rotation ($\Delta x_{\text{COR}} = -11{,}6~\text{pixels}$), du filtrage directionnel des anneaux et du fenêtrage de Ram-Lak a permis d'obtenir un gain en $SNR$ de $+14{,}2~\text{dB}$. 

Pour la seconde séance consacrée aux cubes étalons, cette chaîne de traitement sera appliquée directement. Les défauts calibrés (canaux cylindriques de diamètres $0{,}50 \text{ à } 4{,}00~\text{mm}$ et cavités de sous-extrusion) permettront de quantifier rigoureusement la fonction de transfert de modulation ($MTF$) du système et d'établir la limite de détection en taille de porosité interne pour les thermoplastiques imprimés en 3D.

\section{Conclusion}
\label{sec:conclusion}

Cette première phase d'expérimentation a permis d'établir avec succès une chaîne complète de tomodensitométrie à rayons X pour l'inspection volumique non destructive de pièces issues de la fabrication additive. En mobilisant un banc d'imagerie de laboratoire et un protocole de balayage haute résolution à $720$ projections sur $360^\circ$, les objectifs scientifiques initiaux ont été pleinement atteints.

L'analyse systématique des architectures cinématiques a permis de résoudre l'ambiguïté d'alignement de l'axe de rotation mécanique et de démontrer que la découpe le long des colonnes du capteur (Famille 2) est l'unique configuration fidèle à la géométrie physique du système. La calibration sub-pixel du centre de rotation a révélé un décalage optimal de $\Delta x_{\text{COR}} = -11{,}6 \pm 0{,}2~\text{pixels}$, dont la compensation a éliminé les artéfacts de dédoublement et conféré un gain de rapport signal sur bruit de $+14{,}2~\text{dB}$.

Sur le plan métrologique, l'extraction de profils d'atténuation et la mesure à mi-hauteur (FWHM) ont permis de quantifier l'épaisseur de la coquille extérieure du modèle en PLA à $0{,}81 \pm 0{,}04~\text{mm}$, en accord parfait ($+1{,}25\,\%$, inférieur à l'incertitude à $1\,\sigma$) avec la valeur nominale de consigne ($0{,}80 \pm 0{,}02~\text{mm}$). De surcroît, le réseau interne d'alvéoles de remplissage a été visualisé et mesuré ($3{,}18 \pm 0{,}08~\text{mm}$), confirmant la sensibilité de la tomodensitométrie aux phénomènes de retrait thermique post-extrusion.

Ces résultats valident sans ambiguïté les capacités du dispositif et préparent directement la seconde étape de l'étude. Dès la prochaine séance, ce protocole optimisé sera appliqué aux scans tomographiques de fantômes cubiques étalonnés intégrant des canaux et des cavités de dimensions micrométriques certifiées. Cette suite logique permettra de clore la boucle métrologique en dressant une cartographie quantitative des écarts dimensionnels réels par rapport aux fichiers CAO d'origine.


% -------------------------------------------------------------------------
%                              BIBLIOGRAPHIE
% -------------------------------------------------------------------------
\bibliographystyle{ieeetr}
\bibliography{references}

\end{document}
